\chapter{The Weierstrass Embedding into \texorpdfstring{\(\CP^2\)}{CP2}}
\label{ch:weierstrass-embedding}

The three functions spanning \(L(3O)\) define a projective map
\[
  \Phi:T\longrightarrow\CP^2,\qquad
  z\longmapsto[\wp(z):\wp'(z):1],
\]
with \(O\) sent to \([0:1:0]\).  By \cref{thm:wp-ode-new}, its image lies on
\[
  E:\quad Y^2Z=4X^3-g_2XZ^2-g_3Z^3.
\]

\begin{theorem}
\label{thm:wp-embedding-new}
If \(\Delta=g_2^3-27g_3^2\ne0\), the map \(\Phi\) is a biholomorphism from
\(\C/\Lambda\) onto the smooth projective cubic \(E\).
\end{theorem}

\begin{proof}[Proof sketch]
Near \(O\), the Laurent expansions
\(\wp(z)=z^{-2}+\cdots\) and \(\wp'(z)=-2z^{-3}+\cdots\) show that the map
extends holomorphically to \([0:1:0]\).  The equation shows the image lies on
\(E\).  Smoothness follows from \(\Delta\ne0\).  The evenness of \(\wp\) and
oddness of \(\wp'\) distinguish the two points over a generic \(x\)-value, and
the special half-period points are accounted for by ramification.  Compactness
then promotes the holomorphic bijection to a biholomorphism.
\end{proof}

\begin{keyidea}{The analytic-to-algebraic map}
The map \(\Phi=[\wp:\wp':1]\) is not only a parametrization.  It is the
identification that allows the elementary group law on \(\C/\Lambda\) to become
a geometric law on a cubic.
\end{keyidea}

\section*{Exercises}
\addcontentsline{toc}{section}{Exercises}

\exerciseintent{Compute the projective limit that sends the lattice origin to the point at infinity.}
\begin{exercise}
Use
\[
  \wp(z)=z^{-2}+O(z^2),\qquad \wp'(z)=-2z^{-3}+O(z)
\]
and multiply the homogeneous coordinates
\([\wp(z):\wp'(z):1]\) by \(z^3\).  Compute the limit as \(z\to0\) and obtain
\(\Phi(0)=[0:1:0]\).
\end{exercise}

\exerciseintent{Track how \(z\mapsto -z\) changes the affine Weierstrass coordinates.}
\begin{exercise}
Using \(\wp(-z)=\wp(z)\) and \(\wp'(-z)=-\wp'(z)\), compute
\[
  x(-z),\qquad y(-z)
\]
for \(x(z)=\wp(z)\), \(y(z)=\wp'(z)\).
\end{exercise}

\exerciseintent{Detect singularity of the affine cubic by solving the partial-derivative equations.}
\begin{exercise}
For
\[
  F(x,y)=y^2-4x^3+g_2x+g_3,
\]
compute \(F_x\) and \(F_y\).  Solve the equations
\[
  F=F_x=F_y=0
\]
and show that they are equivalent to \(y=0\) and \(x\) being a repeated root of
\(4x^3-g_2x-g_3\).
\end{exercise}

\exerciseintent{Convert parity of \(\wp,\wp'\) into cubic reflection.}
\begin{exercise}
Compute the homogeneous coordinates of \(\Phi(-z)\) and use the parity of
\(\wp\) and \(\wp'\) to prove
\[
  \Phi(-z)=\ominus\Phi(z),
\]
where \(\ominus(x,y)=(x,-y)\) and \(\ominus O=O\).
\end{exercise}

\exerciseintent{Write one group element in torus, cubic, and divisor-class notation.}
\begin{exercise}
For \(z\in T\), write the same group element in the three notations
\[
  z\in T,\qquad \Phi(z)\in E(\C),\qquad [\Phi(z)]-[O]\in\Pic^0(E).
\]
Give the two explicit formulas
\[
  z\longmapsto[\wp(z):\wp'(z):1],\qquad
  P\longmapsto[P]-[O]
\]
that identify these notations.
\end{exercise}

\exerciseintent{Check the generic injectivity mechanism of the Weierstrass map and mark the theorem boundary.}
\begin{exercise}
Assume \(\Phi(z)=\Phi(w)\) in the affine chart.  Compute the two equalities
\[
  \wp(z)=\wp(w),\qquad \wp'(z)=\wp'(w).
\]
For a generic value of \(\wp\), the equation \(\wp(t)=\wp(z)\) has the two
possibilities \(t\equiv z\) and \(t\equiv -z\) modulo \(\Lambda\).  Use the
oddness of \(\wp'\) to show how the second coordinate removes the generic
\(z\) versus \(-z\) ambiguity.  State which exceptional half-period and
point-at-infinity cases still belong to the full Weierstrass embedding theorem.
\end{exercise}
